\chapter{Giới thiệu chung và cơ sở bài toán}

\section{Bối cảnh tuyển dụng số}

Tuyển dụng số đang dịch chuyển từ mô hình quản lý hồ sơ thụ động sang mô hình ra quyết định dựa trên dữ liệu. Ở mô hình mới, hệ thống không chỉ lưu trữ hồ sơ mà còn phải hỗ trợ đánh giá mức độ phù hợp, giải thích kết quả và theo dõi vòng đời ứng tuyển. Nếu thiếu ba năng lực này, nền tảng tuyển dụng khó tạo giá trị bền vững cho cả ứng viên lẫn nhà tuyển dụng.

Trong thực tế, dữ liệu đầu vào của tuyển dụng có tính nhiễu cao. Hồ sơ cùng một năng lực có thể được mô tả bằng nhiều cách viết khác nhau; mô tả vị trí tuyển dụng cũng thường không đồng nhất về cấu trúc. Do đó, bài toán cốt lõi không nằm ở việc "lưu được dữ liệu", mà ở khả năng chuẩn hóa và liên kết ngữ nghĩa giữa hai phía dữ liệu \cite{requirements}.

\section{Nhóm người dùng và nhu cầu}

Nghiên cứu này xem xét ba nhóm người dùng chính: ứng viên, nhà tuyển dụng và quản trị vận hành. Ứng viên cần công cụ thể hiện năng lực rõ ràng và nhận được gợi ý công việc phù hợp. Nhà tuyển dụng cần cơ chế tìm kiếm, so khớp và lọc hồ sơ hiệu quả, đồng thời có căn cứ để giải thích lựa chọn. Quản trị vận hành cần khả năng theo dõi trạng thái hệ thống, phát hiện sai lệch và đảm bảo quy trình nghiệp vụ được thực thi nhất quán.

\begin{table}[H]
\caption{Nhu cầu cốt lõi theo nhóm người dùng}
\centering
\begin{tabularx}{\textwidth}{p{0.20\textwidth}X}
\toprule
\textbf{Nhóm người dùng} & \textbf{Nhu cầu cốt lõi}\\
\midrule
Ứng viên & Chuẩn hóa hồ sơ cá nhân, nhận gợi ý công việc có giải thích, theo dõi trạng thái ứng tuyển minh bạch.\\
Nhà tuyển dụng & Mô tả vị trí tuyển dụng nhất quán, tìm ứng viên theo tiêu chí bắt buộc và mức độ phù hợp, gửi lời mời có kiểm soát.\\
Quản trị vận hành & Theo dõi chất lượng xử lý tài liệu, giám sát trạng thái quy trình, lưu dấu vết phục vụ kiểm tra và cải tiến.\\
\bottomrule
\end{tabularx}
\end{table}

\section{Phát biểu bài toán}

Bài toán được phát biểu như sau: từ hai tập dữ liệu phi cấu trúc gồm hồ sơ ứng viên và mô tả vị trí tuyển dụng, xây dựng một cơ chế xử lý có khả năng:

\begin{itemize}
  \item Chuẩn hóa dữ liệu thành biểu diễn có cấu trúc.
  \item Áp dụng tiêu chí bắt buộc để loại trừ sớm các cặp không khả thi.
  \item Xếp hạng các cặp còn lại theo mức độ tương đồng ngữ nghĩa.
  \item Trả về kết quả có khả năng giải thích và theo dõi được theo thời gian.
\end{itemize}

Đầu vào của bài toán gồm dữ liệu tài liệu, dữ liệu hồ sơ nghiệp vụ và yêu cầu tìm kiếm/so khớp. Đầu ra của bài toán gồm danh sách gợi ý có thứ tự, điểm đánh giá theo thành phần, nhận xét giải thích và trạng thái nghiệp vụ tương ứng.

\section{Phạm vi MVP và các non-goals}

Phạm vi MVP tập trung vào khả năng vận hành lõi: chuẩn hóa dữ liệu, so khớp hai chiều, hỗ trợ ứng tuyển và lời mời, cùng với cơ chế quan sát vận hành tối thiểu. Những nội dung có độ phức tạp cao hoặc phụ thuộc dữ liệu lớn không nằm trong phạm vi giai đoạn này.

\begin{table}[H]
\caption{Phạm vi nghiên cứu và phần không thuộc phạm vi}
\centering
\begin{tabularx}{\textwidth}{p{0.45\textwidth}p{0.45\textwidth}}
\toprule
\textbf{Trong phạm vi} & \textbf{Ngoài phạm vi}\\
\midrule
Chuẩn hóa hồ sơ và mô tả công việc ở mức MVP. & Tối ưu toàn cục cho hệ sinh thái tuyển dụng quy mô doanh nghiệp lớn.\\
So khớp hai chiều có giải thích. & Kết luận chất lượng xếp hạng ở mức tổng quát khi chưa có dữ liệu gán nhãn lớn.\\
Quy trình ứng tuyển và lời mời có theo dõi trạng thái. & Tự động hóa toàn bộ quyết định nhân sự không có kiểm duyệt của con người.\\
\bottomrule
\end{tabularx}
\end{table}

\section{Các khái niệm nền tảng}

\subsection{Quy trình ATS-lite}

ATS-lite trong nghiên cứu này được hiểu là tập chức năng tinh gọn để quản lý vòng đời tuyển dụng ở mức cần thiết: tạo hồ sơ, tạo nhu cầu tuyển dụng, so khớp, phát sinh tương tác và cập nhật trạng thái. Điểm quan trọng của ATS-lite là giữ được tính kiểm soát quy trình nhưng tránh phức tạp hóa như hệ thống doanh nghiệp toàn diện.

\subsection{Structured parsing}

Structured parsing là bước chuyển đổi dữ liệu văn bản tự do sang dạng có cấu trúc. Mục tiêu của bước này là tạo dữ liệu nhất quán để phục vụ lọc điều kiện và so khớp ngữ nghĩa. Nếu structured parsing thiếu ổn định, toàn bộ pipeline phía sau sẽ suy giảm chất lượng.

\subsection{Embedding và semantic similarity}

Embedding giúp biểu diễn văn bản thành không gian đặc trưng số, cho phép so sánh mức độ gần nhau về nghĩa thay vì chỉ so khớp từ khóa bề mặt. Trong bài toán tuyển dụng, biểu diễn ngữ nghĩa đặc biệt hữu ích với các trường mô tả kinh nghiệm, kỹ năng và yêu cầu vị trí \cite{hld-matching,hld-storage}.

\subsection{Explainable matching}

Explainable matching là nguyên tắc bắt buộc của đề tài: hệ thống không chỉ trả kết quả xếp hạng mà phải cho biết vì sao một kết quả được xếp cao hoặc thấp. Tính giải thích được xây dựng trên hai lớp thông tin: điều kiện bắt buộc và điểm tương đồng theo thành phần.

\section{Mô hình hóa yêu cầu trạng thái}

Để đảm bảo tính nhất quán nghiệp vụ, nghiên cứu sử dụng mô hình trạng thái rõ ràng cho từng thực thể chính. Trạng thái không chỉ dùng để hiển thị mà quyết định khả năng tham gia tìm kiếm và so khớp.

\begin{table}[H]
\caption{Nguyên tắc trạng thái ở mức khái niệm}
\centering
\begin{tabularx}{\textwidth}{p{0.34\textwidth}X}
\toprule
\textbf{Nhóm trạng thái} & \textbf{Nguyên tắc vận hành}\\
\midrule
Hồ sơ ứng viên & Chỉ hồ sơ đã sẵn sàng mới tham gia không gian tìm kiếm công khai.\\
Vị trí tuyển dụng & Chỉ vị trí đang mở mới tham gia so khớp công khai.\\
Quy trình tài liệu & Mỗi lần xử lý tài liệu phải có trạng thái tiến trình và trạng thái lỗi rõ ràng.\\
Tương tác tuyển dụng & Ứng tuyển và lời mời phải có vòng đời trạng thái để truy vết quyết định.\\
\bottomrule
\end{tabularx}
\end{table}

\section{Kết luận chương}

Chương 1 đã xác lập nền tảng bài toán gồm bối cảnh, nhu cầu người dùng, phát biểu kỹ thuật và phạm vi nghiên cứu. Các khái niệm về chuẩn hóa dữ liệu, so khớp ngữ nghĩa và tính giải thích được xem là trục xuyên suốt cho các chương tiếp theo. Trên cơ sở này, chương 2 sẽ phân tích các kỹ thuật phù hợp để hiện thực hóa mô hình đề xuất.
